\clearpage
\fancyhead[L]{\small ASTRONOMISCHES JAHRBUCH.}
\hypertarget{native-616-lower}{}
\begin{center}\small Göttingische gelehrte Anzeigen. Stück 80. Seite $793\;\therefore\;800$. 1819 Mai 20.\\[3pt]\rule{.72\linewidth}{.3pt}\end{center}

\emph{Astronomisches Jahrbuch für das Jahr {\upshape 1821} nebst einer Sammlung der neuesten in die astronomischen Wissenschaften einschlagenden Abhandlungen, Beobachtungen und Nachrichten. Berechnet und herausgegeben von} Dr. J. E. \textsc{Bode}, \emph{Königl. Astronom u. s. w.} Berlin 1818. Bei dem Verfasser und in Commission bei Ferd. Dümmler. 250 Seiten in Octav, nebst einer Kupfertafel.

Da die Einrichtung des astronomischen Kalenders unverändert geblieben ist, so haben wir nur von den begleitenden Aufsätzen Anzeige zu machen. Der erste vom Prof. \textsc{Degen} in Kopenhagen hat das \textsc{Kepler}'sche Problem zum Gegenstande und entwickelt die Coëfficienten der Reihe für die Mittelpunktsgleichung noch auf zwei Ordnungen weiter, als der im vorhergehenden Jahrgange enthaltene Aufsatz von \textsc{Schubert} über dieselbe Aufgabe. Wir beziehen uns dabei auf dasjenige, was wir bei Anzeige des letztern erinnert haben. --- Beobachtete Trabantenverfinsterungen, Sternbedeckungen und Planetenoppositionen vom Prof. \textsc{Bürg} auf der Wiener Sternwarte. --- Beobachtete Solstitien in den Jahren 1812---1817 auf der Turiner Sternwarte von \textsc{Plana}. Auch diese an einem \textsc{Reichenbach}'schen Multiplicationskreise mit stehender Säule, von 18 Zoll Durchmesser, angestellten Beobachtungen geben die Schiefe der Ekliptik aus den Wintersolstitien um 4 Secunden kleiner als aus den Sommersolstitien. Wenn der Herausgeber in einer Anmerkung sagt, dass dies so merkwürdige Phänomen bereits von Astronomen und Physikern in Untersuchung genommen und erklärt sei, und S. 197 noch bestimmter behauptet, dass dasselbe von \textsc{Bessel} \emph{durch eine verbesserte Refractionstafel} beseitigt sei, so beruhet dieses auf einem Irrthum, da jener Unterschied bei Anwendung der \textsc{Bessel}'schen Refractionstafel keinesweges gehoben wird, sondern eher noch etwas grösser ausfällt, als bei Anwendung anderer Refractionstafeln. --- Beobachtete Planetenoppositionen in Kremsmünster von \textsc{Derflinger}. --- Ueber verschiedene astronomische \SourceLeaf{617}Gegenstände von \textsc{Littrow}. Der Kürze wegen schränken wir uns hier auf einige Bemerkungen über die Resultate ein, die Hr. \textsc{Littrow} aus den Königsberger Beobachtungen gezogen hat. Es ist hier das schon früher in der Zeitschrift für Astronomie mitgetheilte Verzeichniss der von Hrn. \textsc{Littrow} aus den Königsberger Beobachtungen abgeleiteten Polardistanzen von 23 Fixsternen wieder abgedruckt. Bekanntlich sind diese Polardistanzen alle um mehrere Secunden grösser, als diejenigen, welche \textsc{Piazzi}, \textsc{Oriani} und \textsc{Pond} mit grosser Uebereinstimmung aus ihren Beobachtungen abgeleitet haben. Es steht zu hoffen, dass in Zukunft die Quelle dieses auffallenden Unterschiedes wird aufgefunden und für die praktische Astronomie lehrreich gemacht werden.. Allein so sehr wir der ausgezeichneten Sorgfalt, womit der vortreffliche Königsberger Astronom die Theilungsfehler seines Kreises bestimmt hat, Gerechtigkeit widerfahren lassen, so scheint uns doch das Urtheil Hrn. \textsc{Littrow}'s, dass die Königsberger Beobachtungen vor allen andern von allen constanten Fehlern gänzlich frei zu sprechen sind, zu voreilig und gewagt. Hr. \textsc{Littrow} theilt hier ferner die von ihm aus den beiden ersten Jahrgängen der Königsberger Beobachtungen abgeleiteten Rectascensionsunterschiede der 36 \textsc{Maskelyne}'schen Fundamentalsterne mit, mit Beifügung der diesen Bestimmungen beizulegenden Gewichte und wahrscheinlichen Fehler. Hätten diese Bestimmungen wirklich eine solche Zuverlässigkeit, wie ihnen hier zugeschrieben wird, so ist kein Zweifel, dass sie alles, was wir in dieser Art bisher besassen, weit übertreffen würden. Allein eben diese unerhörte Genauigkeit, welche man sonst mit den bisherigen Hülfsmitteln nur durch Beobachtungen von einem Menschenalter würde erreichen können, und die angeblich aus zweijährigen Beobachtungen erreicht sein soll, macht eine schärfere Prüfung um so nothwendiger, da dieser Gegenstand für die praktische Astronomie von höchster Wichtigkeit ist. Das neueste Heft der Zeitschrift für Astronomie, worin Hr. \textsc{Littrow} sein Verfahren näher angegeben hat, gibt das Mittel zu dieser Prüfung. Wir sehen daraus, dass Hrn. \textsc{Littrow}'s Behandlung der Beobachtungen für sich schon dem Geiste der Methode der kleinsten Quadrate nicht gemäss, aber was noch viel wichtiger ist, dass seine Bestimmung des Gewichts der Resultate demselben in mehrern wesentlichen Stücken ganz entgegen und dadurch illusorisch ist. Hr. \textsc{Littrow} setzt (Zeitschrift für Astronomie 6. Band S. 23) das Gewicht der Bestimmung des Rectascensionsunterschiedes zweier Sterne, aus den Unterschieden derselben von einem dritten, wenn letztre Bestimmungen die Gewichte $p$ und $\dot{q}$ haben, $=\dfrac{4pq}{p+q}$, welches falsch ist; die richtige Formel ist nur $\dfrac{pq}{p+q}$. Hierdurch allein schon würden die Gewichte seiner Endresultate fast auf ein Viertel ihres angeblichen Werths reducirt werden. Allein noch einflussreicher ist der von Hrn. \textsc{Littrow} übersehene Umstand, dass die Verbindung der partiellen Resultate zu einem Endresultate und die Bestimmung des Gewichts des letztern nur unter der Voraussetzung gültig sind, dass jene alle ganz unabhängig von einander sein müssen, was hier keinesweges statt findet. Hrn. \textsc{Littrow}'s grosse Zahlen werden hiedurch so sehr heruntergebracht, dass z.\,B. die Bestimmung des Rectascensionsunterschiedes von $\gamma$ im Pegasus und $\alpha$ im Schwane, welche angeblich das Gewicht von 1507 einzelnen Vergleichungen haben soll, schwerlich viel mehr als das Gewicht von 40 einzelnen Vergleichungen behalten möchte. Auch gegen einige andere Momente von Hrn. \textsc{Littrow}'s Rechnung liessen sich Erinnerungen machen, die wir hier aber der Kürze wegen übergehen müssen. --- Gegenschein der Pallas 1816, und andere astronomische Beobachtungen 1817 von \textsc{David} und \textsc{Bittner} auf der. Prager Sternwarte. Wir setzen hier nur einige Bemerkungen her über die, an einem 12zolligen Multiplicationskreise von \textsc{Reichenbach} beobachteten, Zenithdistanzen der Sonne und von Fixsternen, dergleichen auch in mehrern der frühern Bände des Jahrbuchs mitgetheilt sind. Es wäre zu wünschen gewesen, dass diese Beobachtungen mehr in einer tabellarischen Form zusammengestellt wären, die die Uebersicht dessen, was sich daraus etwa schliessen lässt, erleichterte. Hr. \textsc{David} hat die Gewohnheit, die beobachtete Zenithdistanz mit derjenigen, die er mit der als bekannt angenommenen Polhöhe und mit der aus Sternverzeichnissen oder den Sonnentafeln \SourceLeaf{618}entlehnten Declination des Gestirns berechnet, zu vergleichen, den Unterschied als die beobachtete Refraction, und die Abweichung dieser von der Tafelrefraction als eine Correction der letztern zu betrachten, ein Verfahren, das wir nicht billigen können. Die Polhöhe von Prag hat Hr. \textsc{David}, so viel wir wissen, nach der \textsc{Hell}'schen Methode bestimmt, durch Sterne, die in Norden und Süden in beinahe gleicher Höhe culminiren, und deren Declinationen er aus den Sternverzeichnissen entlehnt. Das Verfahren ist gut, aber blos für untergeordnete Sternwarten und Instrumente; wer eine so äusserst delicate Sache, wie die Verbesserung der Refractionstafeln, unternimmt, von dem fordert man, dass er seine Polhöhe zuvor selbständig und unabhängig von fremden Bestimmungen festsetze. Was Hr. \textsc{David} als Fehler der Refractionstafeln betrachtet, ist aus der Conspiration des Fehlers der Polhöhe, der Fehler der Declinationen, der Fehler des Instruments, der Fehler der Beobachtungen, Fehler der Aberrations- und übrigen Reductionen und vielleicht zu einem sehr kleinen Theile Fehler der Refractionstafeln entstanden. Diese Fehler von einander zu sondern, dazu reichen nicht einige wenige isolirte Beobachtungen hin; es wird dazu eine Jahrelang mit den besten Instrumenten planmässig ausgeführte und benutzte Reihe von Beobachtungen erfordert. Hr. \textsc{David} glaubt, wenn man auch die absolute Richtigkeit der Abstände selbst bezweifeln wolle, so zeige sich doch in den relativen Verhältnissen der beobachteten Zenithdistanzen der Einfluss der nach den Temperaturänderungen veränderlichen Refractionen. Allein auf diese Aenderungen wird ja schon bei Anwendung der Tafeln selbst Rücksicht genommen; es ist theoretisch bewiesen, dass, wie auch immer die Dichtigkeit der Luftschichten sich ändere, bei so mässigen Zenithdistanzen die Refraction blos von der Dichtigkeit der untersten Luftschicht abhängt, vorausgesetzt,, dass alle Schichten horizontal sind. Dass diese Voraussetzung nicht immer strenge wahr ist, leidet zwar keinen Zweifel, allein der Erfolg \emph{davon} ist nicht eine den Temperaturänderungen folgende, sondern eine unregelmässige keinem Calcül zu unterwerfende Refraction, deren Nachtheile nur durch häufige Beobachtungen und durch Ausschliessung solcher, wo die Gestirne besonders unruhige Bilder gaben, zu heben ist. Wir erkennen daher in den erwähnten Unterschieden nur das Spiel der unvermeidlichen Irregularitäten der Beobachtungen. Noch ist uns aufgefallen, dass Hr. \textsc{David} die mit dem Crystallwürfel (Prisma) beobachteten Zenithdistanzen immer um $1''$ vermindert; wir können diese angebliche Correction nicht billigen, da wir in der Natur des Instruments keinen Grund finden, der zur Voraussetzung eines für alle Zenithdistanzen constanten Fehlers berechtigte. --- Noch von derselben Sternwarte Beobachtungen des Jupiter, Uranus, Saturn und der Ceres. --- Comet vom 1. Nov. 1817, welcher vom Dr. \textsc{Olbers} entdeckt, aber nur ein einzigesmal beobachtet wurde. --- Beobachtungen des (ersten) Cometen von 1818 vom Dr. \textsc{Olbers}. --- Ueber einige merkwürdige Stellen der Milchstrasse von Dr. \textsc{Herschel} (Auszug aus den Philosoph. Transact.). --- Beobachtungen von Planetenoppositionen u. s. w. in Wilna von \textsc{Sniadecki}. --- Beobachtungen und Bestimmungen der Bahn des ersten Cometen von 1818, von \textsc{Enke}. Merkwürdig ist an diesem Cometen seine Lichtschwäche, welche zunahm, als sein Licht (mochte man es als eignes oder erborgtes betrachten) nach der Rechnung noch hätte wachsen sollen; wahrscheinlich erlitt der Comet während der Dauer seiner Sichtbarkeit bedeutende innere Veränderungen. --- Beobachtete Trabantenverfinsterungen, Sternbedeckungen und Sonnenflecke im Jahre 1817, von \textsc{Hallashka} in Prag. --- Formeln zur Berechnung des Orts eines Gestirns, aus beobachteten Allignements mit vier Sternen von \textsc{Bessel}. Sie beruhen, dem Wesentlichen nach, auf der vorgängigen Bestimmung der Durchschnittspunkte der beiden grössten Kreise, in welchen die zwei bekannten Sternpaare liegen, mit dem Aequator (oder der Ekliptik), obwohl die zierliche Einkleidung der Auflösung rein analytisch ist. --- Von demselben Astronomen Beobachtungen von Planetenoppositionen, Sternbedeckungen, Solstitien und dem Polarstern im J. 1817. --- Auszug aus den astronomischen Beobachtungen auf der Berliner Sternwarte in demselben Jahre. --- Bestimmung der Polhöhe von Krakau. Mit \SourceLeaf{619}einem \textsc{Reichenbach}'schen multiplicirenden Theodolithen fand Hr. Kreiscommissarius \textsc{Lorenz} aus Beobachtungen des Nordsterns $50^\circ\,3'\,57''7$, aus Sonnenbeobachtungen $50^\circ\,3'\,50''$; mit dem Mittel stimmt sehr nahe Hrn. \textsc{Sniadecky}'s ältere Bestimmung überein, die später, und, wie aus den neuen Beobachtungen hervorzugehen scheint, mit Unrecht, durch \textsc{Littrow} um $17''$ vermindert war. --- Ueber die Genauigkeit der astronomischen Beobachtungen und Berechnungen. Dieser Aufsatz hat zum Zweck, anschaulich zu machen, wie sehr viel es heutzutage schwerer ist, in der Astronomie neue Fortschritte zu thun, als vormals, wo die Forderungen an den beobachtenden und rechnenden Astronomen so viel geringer waren. Man möchte, sagt der Verf., die Lage der ältern Astronomen beneiden. Sie konnten mit weit weniger Anstrengung den Himmelslauf beobachten und berechnen, da noch keine grosse oder die grösste Genauigkeit erwartet werden konnte, und wurden nichts destoweniger bei ihren Zeitgenossen und bei der Nachwelt berühmt u.\,s.\,w. So wahr das ist, wenn wir die Leistungen der ältern Astronomen aus \emph{unserm} Standpunkte betrachten, so darf man doch auch nicht aus der Acht lassen, wie unbehülflich und unvollkommen ihre Instrumente, wie beschwerlich ihre Beobachtungsmethoden, wie unausgebildet der Calcul war. Nimmt man darauf billige Rücksicht, so möchte wohl z.\,B. ein \textsc{Tycho} seine verdiente Celebrität als praktischer Astronom nicht wohlfeiler gehabt haben, als ein \textsc{Piazzi}. Ueber einzelne Aeusserungen dieses Aufsatzes, die heutige praktische Astronomie betreffend, liesse sich doch noch manches erinnern. Nachdem z.\,B. der Verf. berechnet hat, dass bei einem Kreise von 12 Zoll Halbmesser ein Strich, den hundertsten Theil einer Linie dick, schon über 14 Secunden deckt, fragt er: Sollte eine Menschenhand eine Duodecimallinie noch in 100 kenntliche Striche oder Punkte einzutheilen im Stande sein? Allein die wirkliche Eintheilung der Decimallinie durch 100 Striche wird ja nicht verlangt: dass aber die Eintheilung in diejenige Anzahl von Theilen, die zweckmässig ist, mit einer solchen Genauigkeit von den ersten Künstlern ausgeführt werden könne, dass die Fehler \emph{viel} weniger als den hundertsten Theil einer Linie erreichen, beweisen die Meisterwerke eines \textsc{Reichenbach}, \textsc{Repsold}, \textsc{Troughton} \emph{durch die That}. --- Neue Elemente der Junobahn, Beobachtungen dieses Planeten und der Sonne und andere astronomische Nachrichten von \textsc{Nicolai}; in den letztern wird des \textsc{Repsold}'schen Kreises auf der hiesigen Sternwarte erwähnt, wo aber der vielleicht durch einen Druckfehler zu 5 Fuss angegebene Durchmesser dieses Instruments auf $3\tfrac{1}{2}$ Pariser Fuss zu berichtigen ist. --- Beobachtung der Mondfinsterniss vom 20. April 1818 in Dresden. --- Gerade Aufsteigungen von 35 der vornehmsten Sterne für 1818, nach \textsc{Maskelyne}'s letzten Beobachtungen, aus dem Nautical Almanac für 1820. (Wir bemerken hiebei, dass im Jahrgange des Nautical Almanac für 1821 ein noch ausgedehnteres und auf \textsc{Pond}'s Beobachtungen mit dem neuen Mittagsfernrohr gegründetes Verzeichniss sich findet.) Beobachtungen der scheinbaren Grösse verschiedner Sterne, von den Jahren 1704---1709 (aus des Berliner Astronomen \textsc{Gottfr. Kirch} nachgelassenen Papieren). --- Ueber den \textsc{Repsold}'schen Meridiankreis, nebst verschiedenen astronomischen Beobachtungen von \textsc{Gauss}. --- Beobachtungen der Sonnenfinsterniss vom 5. Mai 1818 von \textsc{Struve}. --- Beobachtung der Opposition der Vesta im Jahr 1818 von \textsc{Enke}. --- Dann ferner Ephemeriden für die Vesta, Juno und Pallas von \textsc{Enke}, \textsc{Posselt}, \textsc{Nicolai} und \textsc{Dirksen}. --- Die übrigen kleinern Artikel und Notizen hier einzeln anzuführen, verstattet der Raum nicht.
